\documentclass{article}
\usepackage[T1]{fontenc}
\usepackage{lmodern}
\usepackage{graphicx}
\usepackage{amsmath,amssymb}
\usepackage[a4paper,margin=2cm]{geometry}
\usepackage[absolute,overlay]{textpos}
\setlength{\TPHorizModule}{1mm}
\setlength{\TPVertModule}{1mm}
\usepackage[indonesian]{babel}
\usepackage{hyperref}
\setlength{\emergencystretch}{2em}
% unit: Halaman 22

\begin{document}

% === Halaman 22 (python/native) ===

Perkalian Titik

• Perkalian titik:  kP = Q   

 Ket:  k adalah skalar, P dan Q adalah titik pada kurva eliptik

• Perkalian titik diperoleh dengan perulangan dua operasi dasar kurva eliptik 

yang sudah dijelaskan:

 1. Penjumlahan titik (P + Q = R)

 2. Penggandaan titik (2P = R)

• Contoh: k = 3 → 3P = P + P + P atau 3P = 2P + P 

    k = 23 → kP = 23P = 2(2(2(2P) + P) + P) + P

22

\end{document}
